\documentclass[12pt,a4paper]{article}
\usepackage[utf8]{inputenc}
\usepackage[spanish,es-tabla]{babel}
\usepackage{geometry}
\usepackage{graphicx}
\usepackage{multirow}
\usepackage{tabularx}
\usepackage{mathptmx}
\usepackage{fancyhdr}
\usepackage{array}
\usepackage[table]{xcolor}
\usepackage{titlesec}
\usepackage{setspace}
\usepackage{microtype}
\usepackage{lastpage}
\usepackage{booktabs}
\usepackage{enumitem}
\usepackage{amsmath}
\usepackage{amsfonts}
\usepackage{amssymb}

\usepackage{float}
\usepackage{placeins}
\usepackage{tikz}
\usepackage{pgfplots}
\pgfplotsset{compat=1.17}
\usetikzlibrary{arrows.meta,positioning,shapes.geometric,calc}

\graphicspath{{images/}{../images/}{./}}

\geometry{
    top=4.7cm,
    bottom=2.5cm,
    left=2.5cm,
    right=2.5cm,
    headheight=3.9cm,
    headsep=0.4cm
}

\setstretch{1.15}
\setlength{\parindent}{0pt}
\setlength{\parskip}{0.45em}
\setlength{\tabcolsep}{4pt}

\titleformat{\section}
  {\normalfont\fontsize{12}{14.5}\bfseries}{\thesection.}{0.6em}{\uppercase}
\titlespacing*{\section}
  {0pt}{1.8ex plus 0.4ex minus 0.2ex}{0.6ex plus 0.2ex}

\titleformat{\subsection}
  {\normalfont\fontsize{11.5}{13.5}\bfseries}{\thesubsection.}{0.5em}{}
\titlespacing*{\subsection}
  {0pt}{1.4ex plus 0.3ex minus 0.1ex}{0.4ex plus 0.1ex}

\titleformat{\subsubsection}
  {\normalfont\fontsize{10.5}{12.5}\bfseries\itshape}{\thesubsubsection.}{0.4em}{}
\titlespacing*{\subsubsection}
  {0pt}{1.1ex plus 0.2ex minus 0.1ex}{0.3ex plus 0.1ex}

\newcolumntype{Y}{>{\centering\arraybackslash}X}
\newcolumntype{L}{>{\raggedright\arraybackslash}X}
\newcolumntype{C}{>{\centering\arraybackslash}X}

\definecolor{headerblue}{RGB}{26, 54, 93}
\definecolor{darkblue}{RGB}{43, 108, 176}
\definecolor{lightgraybg}{RGB}{248, 248, 248}
\definecolor{tableheader}{RGB}{232, 238, 246}
\definecolor{navyheader}{RGB}{0, 32, 96}
\definecolor{lineblue}{RGB}{0, 102, 204}
\definecolor{areablue}{RGB}{220, 236, 252}
\definecolor{lcsred}{RGB}{217, 56, 41}
\definecolor{lcigreen}{RGB}{46, 125, 50}
\definecolor{lciorange}{RGB}{230, 81, 0}
\definecolor{ishikawared}{RGB}{160, 0, 0}

\tikzset{
    catbox/.style={draw, fill=tableheader, font=\footnotesize\bfseries, minimum width=2.8cm, minimum height=0.65cm, align=center},
    efectobox/.style={draw=ishikawared!90!black, line width=1.2pt, fill=ishikawared, text=white, font=\small\bfseries, align=center, inner sep=6pt, text width=3.8cm},
    causetext/.style={font=\scriptsize, align=right, text=black!90}
}

\newcommand{\docCodigo}{TI-SIGE-ENT-F1-006}
\newcommand{\docVersion}{1.0}
\newcommand{\docProceso}{EVALUACIÓN FINANCIERA / PROJECT CHARTER}
\newcommand{\docElaboracion}{24/08/2026}
\newcommand{\docRevision}{24/08/2026}
\newcommand{\docTitulo}{ELABORACIÓN DE CASO DE NEGOCIO, ROI Y PROJECT CHARTER}
\newcommand{\docOrganizacion}{FACULTAD DE INGENIERÍA EN SISTEMAS, ELECTRÓNICA E INDUSTRIAL}
\newcommand{\docSubOrganizacion}{CARRERA DE INGENIERÍA DE SOFTWARE}

\newcommand{\encabezadofrecuente}{%
    \renewcommand{\arraystretch}{1.15}%
    \noindent\makebox[\textwidth][c]{%
    \begin{tabularx}{\textwidth}{|>{\centering\arraybackslash}m{2.8cm}|Y|>{\centering\arraybackslash}m{4.8cm}|}
        \hline
        \multirow{4}{2.8cm}{\centering\raisebox{-1.15cm}{\includegraphics[width=2.5cm,height=2.0cm,keepaspectratio]{images/logo.png}}}
        & \textbf{\fontsize{9}{11}\selectfont \docOrganizacion} 
        & \textbf{\fontsize{8}{10}\selectfont CÓDIGO:} \fontsize{8}{10}\selectfont \docCodigo \\
        \cline{3-3}
        & \textbf{\fontsize{8.5}{10.5}\selectfont \docSubOrganizacion} 
        & \textbf{\fontsize{8}{10}\selectfont VERSIÓN:} \fontsize{8}{10}\selectfont \docVersion \\
        \cline{2-3}
        & \textbf{\fontsize{8}{10}\selectfont PROCESO:} \fontsize{8}{10}\selectfont \docProceso 
        & \textbf{\fontsize{8}{10}\selectfont ELABORACIÓN:} \fontsize{8}{10}\selectfont \docElaboracion \\
        \cline{2-3}
        & \textbf{\fontsize{8.5}{10.5}\selectfont \docTitulo} 
        & \textbf{\fontsize{8}{10}\selectfont ÚLTIMA REVISIÓN:} \fontsize{8}{10}\selectfont \docRevision \\
        \hline
    \end{tabularx}%
    }%
}

\pagestyle{fancy}
\fancyhf{}
\fancyhead[C]{\encabezadofrecuente}
\fancyfoot[C]{\fontsize{10}{12}\selectfont Página \thepage\ de \pageref{LastPage}}
\renewcommand{\headrulewidth}{0pt}

\begin{document}

\noindent{\fontsize{10.5}{12.5}\selectfont \textbf{Integrantes:} Heidi Villavicencio \quad$\vert$\quad Steeven Toala \quad$\vert$\quad Nixon Hurtado \quad$\vert$\quad Andrea Vásquez}

\vspace{0.25cm}

\section{Título}
\textbf{Denominación Oficial del Entregable:} Elaboración de Caso de Negocio, ROI y Project Charter -- Código EDT: 1.6

\section{Introducción}
Todo proyecto de ingeniería de software en el sector público universitario debe fundamentarse en una estricta evaluación de viabilidad económica, operativa y estratégica. El presente informe formaliza el Caso de Negocio del sistema SIGE, demostrando cuantitativamente la rentabilidad institucional mediante el Retorno de Inversión (ROI), el Valor Actual Neto (VAN), la Tasa Interna de Retorno (TIR) y el Acta de Constitución formal del proyecto (\textit{Project Charter}) con la asignación presupuestaria de \$630.30 para gastos de infraestructura local.

\section{Objetivo}
\subsection{Objetivo General}
Evaluar y fundamentar la viabilidad económica, operativa y financiera del sistema SIGE mediante el cálculo de indicadores de rentabilidad (ROI, VAN, TIR) y la emisión formal del Project Charter del proyecto.
\subsection{Objetivos Específicos}
\begin{itemize}[leftmargin=1.5em, itemsep=2pt]
    \item Cuantificar los costos operativos del proceso actual manual frente a los ahorros anuales proyectados con SIGE.
    \item Calcular los indicadores financieros de rentabilidad (ROI $> 300\%$, Payback $< 6$ meses).
    \item Emitir el Project Charter que autoriza formalmente al equipo de ingeniería el uso de los recursos institucionales.
\end{itemize}

\section{Alcance}
\subsection{Límites del Entregable (In-Scope)}
\begin{itemize}[leftmargin=1.5em, itemsep=2pt]
    \item Evaluación financiera a horizonte de 4 años para 65 docentes de la FISEI.
    \item Estructura de costos de infraestructura local (Docker, MinIO y servidor institucional).
    \item Acta de Constitución formal con roles del equipo.
\end{itemize}
\subsection{Exclusiones (Out-of-Scope)}
\begin{itemize}[leftmargin=1.5em, itemsep=2pt]
    \item Adquisición de servidores físicos nuevos (se reutiliza infraestructura existente de la facultad).
\end{itemize}

\section{Definiciones}
\begin{itemize}[leftmargin=1.5em, itemsep=2pt]
    \item \textbf{Project Charter (Acta de Constitución):} Documento emitido por el patrocinador que autoriza formalmente la existencia del proyecto.
    \item \textbf{ROI (Return on Investment):} Relación porcentual entre los beneficios netos obtenidos y los costos incurridos en la inversión.
    \item \textbf{VAN (Valor Actual Neto):} Suma del valor presente de los flujos de caja futuros descontados a una tasa de oportunidad institucional.
    \item \textbf{Payback:} Tiempo necesario para que los ahorros acumulados amortigüen la inversión inicial.
\end{itemize}

\section{Responsabilidades}
\begin{itemize}[leftmargin=1.5em, itemsep=2pt]
    \item \textbf{Equipo de Proyecto SIGE (ST, HV, NH, AV - EDT 1.6):} Modelado financiero, estimación de costos y redacción del Project Charter.
    \item \textbf{Ing. Paulo Torres (Docente Evaluador / Patrocinador):} Aprobación del caso de negocio y suscripción del Project Charter.
\end{itemize}

\section{Normativa Legal y Técnica}
\subsection{Marco Legal Ecuatoriano}
\begin{itemize}[leftmargin=1.5em, itemsep=2pt]
    \item \textbf{Ley Orgánica de Educación Superior (LOES):} Eficiencia en la utilización de los recursos públicos de las universidades.
    \item \textbf{Código Orgánico de Planificación y Finanzas Públicas:} Criterios de evaluación social y beneficio institucional.
\end{itemize}

\subsection{Normas Técnicas Internacionales}
\begin{itemize}[leftmargin=1.5em, itemsep=2pt]
    \item \textbf{Guía PMBOK (7ma Edición), Inicio del Proyecto:} Desarrollo del Acta de Constitución del Proyecto.
    \item \textbf{IEEE Std 1058-1998:} Estándar para planes de gestión de proyectos de software.
\end{itemize}

\section{Método, Políticas y Contenidos}
\subsection{Presupuesto y Evaluación Financiera}
El presupuesto asignado se calculó considerando licencias de desarrollo, dominio y contingencias operativas:

\begin{table}[H]
\centering
\small
\renewcommand{\arraystretch}{1.2}
\begin{tabularx}{\textwidth}{|c|>{\bfseries}m{4.5cm}|X|c|}
\hline
\rowcolor{tableheader}
\textbf{Ítem} & \textbf{Rubro de Inversión} & \textbf{Justificación Técnica} & \textbf{Monto USD} \\
\hline
1 & Infraestructura de Contenedores & Despliegue en VPS local con Docker y MinIO & \$0.00 (FOSS) \\
\hline
2 & Base de Datos PostgreSQL 16 & Motor relacional Open Source sin costo de licenciamiento & \$0.00 (FOSS) \\
\hline
3 & Certificado SSL y Dominio & Seguridad HTTPS TLS v1.3 y nombre de dominio institucional & \$38.00 \\
\hline
4 & Asistente IA (Gemini API) & Tokens para capa de sugerencias ortográficas y síntesis & \$50.00 \\
\hline
5 & Reserva de Imprevistos & Fondo de contingencia operativa ($10\%$) & \$42.30 \\
\hline
6 & Horas de Ingeniería & Trabajo de desarrollo y validación académica del equipo & \$500.00 \\
\hline
\multicolumn{3}{|r|}{\textbf{PRESUPUESTO TOTAL INICIAL ASIGNADO:}} & \textbf{\$630.30} \\
\hline
\end{tabularx}
\caption{Presupuesto Consolidado de Inversión del Proyecto SIGE.}
\label{tab:presupuesto}
\end{table}
\FloatBarrier

\section{Descripción de actividades}
\subsection{Cuantificación del Ahorro Anual y Retorno de Inversión (ROI)}
Considerando una planta de 65 docentes que elaboran 2 informes semestrales al año (130 informes anuales), con un costo por hora docente estimado de \$15.00:
\begin{itemize}[leftmargin=1.5em, itemsep=2pt]
    \item \textbf{Horas Invertidas en Proceso Manual:} $130 \times 8 \text{ h} = 1,040 \text{ h/año}$ (Costo: \$15,600).
    \item \textbf{Horas Invertidas en Revisiones y Reprocesos:} $600 \text{ h/año}$ de coordinadores (Costo: \$9,000).
    \item \textbf{Horas Requeridas con SIGE:} $130 \times 1.5 \text{ h} = 195 \text{ h/año}$ (Costo: \$2,925).
    \item \textbf{Ahorro Neto Anual Estimado:} \textbf{\$21,675.00 anuales}.
\end{itemize}

\begin{table}[H]
\centering
\small
\renewcommand{\arraystretch}{1.2}
\begin{tabularx}{\textwidth}{|>{\bfseries}m{4.5cm}|c|X|}
\hline
\rowcolor{tableheader}
\textbf{Indicador Financiero} & \textbf{Valor Calculado} & \textbf{Interpretación Técnica y Viabilidad} \\
\hline
Retorno de Inversión (ROI) & \textbf{+3,438\%} & Rendimiento extremadamente alto por automatización de procesos burocráticos. \\
\hline
Valor Actual Neto (VAN a 12\%) & \textbf{\$64,820.00} & Proyecto altamente viable que crea valor social y académico directo. \\
\hline
Tasa Interna de Retorno (TIR) & \textbf{$>180\%$} & Supera ampliamente la tasa de corte institucional. \\
\hline
Período de Recuperación (Payback) & \textbf{1.8 meses} & La inversión inicial se amortiza antes del cierre del primer semestre lectivo. \\
\hline
\end{tabularx}
\caption{Indicadores Financieros de Evaluación del Proyecto SIGE.}
\label{tab:roi}
\end{table}
\FloatBarrier

\subsection{Acta de Constitución Oficial (Project Charter)}
Se formaliza el nombramiento del equipo de desarrollo integrado por Steeven Toala (Líder / Arquitecto), Heidi Villavicencio (Analista de Calidad / UI-UX), Nixon Hurtado (Backend / DBA) y Andrea Vásquez (Frontend / QA), bajo la dirección y patrocinio institucional del Ing. Paulo Torres.

\section{Conclusiones}
\begin{itemize}[leftmargin=1.5em, itemsep=2pt]
    \item El caso de negocio demuestra la altísima rentabilidad social y financiera del proyecto, amortizándose en menos de 2 meses de operación.
    \item La suscripción del Project Charter formaliza la existencia del proyecto y blinda el presupuesto asignado de \$630.30.
\end{itemize}

\section{Firmas de Responsabilidad}

\noindent
\renewcommand{\arraystretch}{1.2}
\begin{tabularx}{\textwidth}{|>{\raggedright\arraybackslash}m{2.9cm}|>{\raggedright\arraybackslash}X|>{\centering\arraybackslash}m{4.6cm}|>{\centering\arraybackslash}m{2.3cm}|}
    \hline
    \rowcolor{tableheader}
    \textbf{ACCIÓN} & \textbf{NOMBRE Y CARGO} & \textbf{FIRMA} & \textbf{FECHA} \\
    \hline
    \textbf{Elaborado por:} & Heidi Villavicencio\newline Analista de Calidad / UI-UX & \parbox[c]{4.4cm}{\centering\vspace{4pt}\includegraphics[height=1.5cm,width=4.0cm,keepaspectratio]{images/firma_heidi.png}\vspace{4pt}} & 24/08/2026 \\
    \hline
    \textbf{Revisado por:} & Steeven Toala\newline Líder de Proyecto / Arquitecto & \parbox[c]{4.4cm}{\centering\vspace{4pt}\includegraphics[height=1.5cm,width=4.0cm,keepaspectratio]{images/firma_steveen.png}\vspace{4pt}} & 24/08/2026 \\
    \hline
    \textbf{Aprobado por:} & Ing. Paulo Torres, Mg.\newline Docente Evaluador -- Gestión Proyectos & \parbox[c]{4.4cm}{\centering\vspace{8pt}\textbf{\textcolor{gray!80}{\scriptsize [ PENDIENTE DE APROBACIÓN ]}}\vspace{8pt}} & \textbf{\textcolor{gray!90}{Pendiente}} \\
    \hline
\end{tabularx}

\vspace{0.4cm}

\section{Control de Cambios}

\noindent
\renewcommand{\arraystretch}{1.2}
\begin{tabularx}{\textwidth}{|>{\hsize=0.45\hsize\centering\arraybackslash}X|>{\hsize=1.55\hsize\raggedright\arraybackslash}X|>{\hsize=1.2\hsize\raggedright\arraybackslash}X|>{\hsize=0.8\hsize\centering\arraybackslash}X|}
    \hline
    \rowcolor{tableheader}
    \textbf{Versión} & \textbf{Descripción del cambio} & \textbf{Responsable del cambio} & \textbf{Fecha de aprobación} \\
    \hline
    1.0 & Emisión inicial del entregable según EDT del Tablero Trello & Equipo de Proyecto SIGE & 24/08/2026 \\
    \hline
\end{tabularx}

\vspace{0.4cm}

\section{Anexos}

\subsection*{Anexo 1: Ficha de Registro y Control en el Tablero Oficial de Gestión (EDT Trello)}
\begin{table}[H]
\centering
\small
\renewcommand{\arraystretch}{1.2}
\begin{tabularx}{\textwidth}{|>{\bfseries}m{3.5cm}|X|}
\hline
\rowcolor{tableheader}
\textbf{Parámetro EDT} & \textbf{Detalle Oficial del Tablero Trello} \\
\hline
Código de Tarea & EDT 1.6 \\
\hline
Nombre de la Tarjeta & Elaboración de Caso de Negocio, ROI y Project Charter \\
\hline
Fase Correspondiente & FASE 1. ANÁLISIS Y PLANIFICACIÓN \\
\hline
Período de Ejecución & 20 de agosto al 24 de agosto de 2026 (3 días laborables) \\
\hline
Horas Asignadas & 24 horas hombre netas \\
\hline
Responsables Técnicos & Todo el Equipo (ST, HV, NH, AV) \\
\hline
Entregable Oficial & \texttt{06\_Caso\_de\_Negocio\_ROI\_Project\_Charter.pdf} \\
\hline
Estado en el Tablero & 100\% Completado y Verificado en Reunión de Hito 1 \\
\hline
\end{tabularx}
\caption{Ficha técnica de la tarea EDT 1.6 en el Tablero Trello.}
\end{table}

\subsection*{Anexo 2: Flujo de Caja Descontado Proyectado a 4 Años}
\begin{table}[H]
\centering
\footnotesize
\renewcommand{\arraystretch}{1.15}
\begin{tabularx}{\textwidth}{|c|c|c|c|c|}
\hline
\rowcolor{tableheader}
\textbf{Año} & \textbf{Inversión / Costo Operativo} & \textbf{Ahorro Bruto} & \textbf{Flujo Neto} & \textbf{Flujo Descontado (12\%)} \\
\hline
Año 0 & \$630.30 & \$0.00 & -\$630.30 & -\$630.30 \\
\hline
Año 1 & \$120.00 & \$21,675.00 & \$21,555.00 & \$19,245.54 \\
\hline
Año 2 & \$120.00 & \$21,675.00 & \$21,555.00 & \$17,183.51 \\
\hline
Año 3 & \$120.00 & \$21,675.00 & \$21,555.00 & \$15,342.42 \\
\hline
Año 4 & \$120.00 & \$21,675.00 & \$21,555.00 & \$13,698.59 \\
\hline
\multicolumn{4}{|r|}{\textbf{VALOR ACTUAL NETO (VAN):}} & \textbf{\$64,839.76} \\
\hline
\end{tabularx}
\caption{Flujo financiero descontado del proyecto SIGE.}
\end{table}

\subsection*{Anexo 3: Trazabilidad en Repositorio Institucional Git}
\begin{itemize}[leftmargin=1.5em, itemsep=2pt]
    \item \textbf{Repositorio:} \texttt{FISEI\_Project / SIGE}
    \item \textbf{Rama de Trabajo:} \texttt{main}
    \item \textbf{Commit de Consolidación:} \texttt{e12480a - docs(finance): modelo financiero de retorno de inversion ROI y project charter}
    \item \textbf{Archivos Afectados:} \texttt{docs/caso\_negocio\_roi.md}, \texttt{Entregables\_Latex/Fase\_1/1.6*.tex}
\end{itemize}


\end{document}
